%% Related Work section

\section{Related Work}
\label{sec:related_work}


\textbf{On-chain Security and Program Analysis.}
The immutability of deployed smart contracts has driven extensive research on on-chain security analysis, broadly encompassing pre-deployment verification and post-deployment transaction monitoring. Pre-deployment approaches rely on static analysis and formal verification, with symbolic execution tools such as Oyente~\cite{Luu2016Making} and Maian~\cite{Nikolic2018Finding}, and mature frameworks like Slither~\cite{slither} and Mythril~\cite{mythril2018}, enabling automated detection of common vulnerabilities. Formal methods, including Certora Prover~\cite{certoraProver} and KEVM~\cite{Hildenbrandt2018KEVM}, further strengthen correctness guarantees by reasoning over mathematically precise semantics. However, static analysis is fundamentally constrained by path explosion and its inability to capture the non-deterministic, evolving global state of blockchain systems, particularly in DeFi environments characterized by composability~\cite{defi2023} and gas-related vulnerabilities~\cite{MadMax}. These limitations have motivated post-deployment transaction analysis, where graph-based methods, machine learning techniques~\cite{mythischaikul2023machine}, and forensic tools such as TXSPECTOR~\cite{TXSPECTOR} are used to identify malicious behaviors including Ponzi schemes and flash-loan-based attacks~\cite{zhou2021flashboys}. Nevertheless, existing dynamic approaches largely depend on predefined heuristics, leaving open the challenge of detecting malicious interactions that remain syntactically compliant yet exploit subtle business logic manipulations.

%\subsection{On-chain Security and Program Analysis}
% \textbf{On-chain Security and Program Analysis.}
% The escalating complexity of decentralized ecosystems has rendered the security posture of smart contracts increasingly precarious. Given the immutable nature of blockchain code once deployed, the defensive paradigm has evolved from rudimentary static audits toward sophisticated, real-time transaction analysis. Current methodologies generally bifurcate into two distinct categories: proactive pre-deployment verification and retrospective runtime monitoring.

% The former primarily leverages pattern-based static analysis and formal verification. Early symbolic execution tools like Oyente \cite{Luu2016Making} and Maian \cite{Nikolic2018Finding} laid the groundwork for automated vulnerability detection. Modern foundational frameworks such as Slither \cite{slither} and Mythril \cite{mythril2018} utilize symbolic execution and taint analysis to identify classic vulnerabilities like access control flaws. To achieve higher reliability, frameworks like Certora Prover \cite{certoraProver} translate code into mathematical models to prove security invariants, while KEVM \cite{Hildenbrandt2018KEVM} provides a complete formal semantics of the Ethereum Virtual Machine for rigorous verification.

% Despite their utility, these static approaches are often bottlenecked by path explosion in complex loops and a fundamental inability to simulate the non-deterministic global state of the blockchain. Such limitations—particularly the difficulty in accounting for DeFi composability \cite{defi2023} and gas-related vulnerabilities like those identified by MadMax \cite{MadMax}—have catalyzed the growth of post-deployment transaction analysis. While existing research has successfully applied graph analysis, machine learning \cite{mythischaikul2023machine}, and specialized transaction forensic tools like TXSPECTOR \cite{TXSPECTOR} to characterize malicious behaviors (e.g., Ponzi schemes or flash loan manipulations \cite{zhou2021flashboys}), these dynamic methods remain heavily reliant on pre-defined heuristics. This raises a critical question: how can we detect ``seemingly compliant'' malicious interactions that bypass syntactic rules through sophisticated business logic manipulation?

%\subsection{LLMs for Code and Security Analysis}
\noindent
\textbf{LLMs for Code and Security Analysis.}
The emergence of Large Language Models (LLMs) has initiated a paradigm shift in smart contract auditing, steering the field away from rigid pattern matching toward nuanced semantic inference. Early breakthroughs were rooted in representation learning, where models like CodeBERT \cite{feng-etal-2020-codebert} and SmartBERT demonstrated a profound capacity for capturing program control flows and semantics through domain-specific pre-training. Building upon this foundation, the focus has shifted toward fine-grained defect localization, with models such as VulBERTa \cite{hanif2022vulberta} and LineVul \cite{fu2022linevul} significantly enhancing the precision of vulnerability prediction.

More recently, LLM-driven frameworks like GPTScan \cite{GPTScan} have begun integrating reasoning capabilities—often enhanced by Chain-of-Thought (CoT) prompting \cite{wei2023chainofthoughtpromptingelicitsreasoning}—with program analysis to validate potential attack vectors. Nevertheless, deploying LLMs in the high-stakes DeFi environment is not without its perils. The most pervasive challenge is the phenomenon of ``hallucinations'', wherein models generate plausible yet technically erroneous vulnerability reports \cite{10.1145/3442188.3445922}. Furthermore, the inherent ``knowledge cutoff'' of pre-trained models renders them blind to zero-day exploits and rapidly evolving protocol patterns. This lack of interpretability, often referred to as the ``black-box'' nature of LLMs \cite{khare2024understandingeffectivenesslargelanguage}, and their performance degradation when operating outside their ``comfort zone'' \cite{guo2024outsidecomfortzoneanalysing}, further complicates their adoption by security professionals who require verifiable evidence for decision-making.

%\subsection{RAG in Software Engineering and Security}
\noindent
\textbf{RAG in Software Engineering and Security.}
To mitigate the inherent limitations of static parameters, Retrieval-Augmented Generation (RAG) has surfaced as a cornerstone for building knowledge-intensive security systems \cite{lewis2021retrievalaugmentedgenerationknowledgeintensivenlp}. By bridging the gap between a model's parametric memory and non-parametric external databases, RAG effectively anchors LLM outputs in real-world facts, thereby enhancing both temporal relevance and factual accuracy.

While the efficacy of RAG has been validated in broader cybersecurity tasks—ranging from repository-level code completion in RepoCoder \cite{zhang2023repocoder} and documentation-driven generation in DocPrompting \cite{zhou2023docpromptinggeneratingcoderetrieving} to LLM-empowered penetration testing in PentestGPT \cite{deng2024pentestgptllmempoweredautomaticpenetration}—its application within the smart contract domain remains in its infancy. Preliminary efforts have begun utilizing retrieval mechanisms to assist in generating security invariants, such as SmartInv \cite{wang2024smartinvmultimodallearningsmart}; however, a significant research gap persists: the absence of a dynamic, temporal knowledge base designed for real-time attack detection. Conventional RAG implementations are largely static, prioritizing isolated code snippets over the multi-dimensional context of live transactions, such as state transitions and economic repercussions.

This study addresses this gap by introducing a Semantically-Enhanced Temporal Transaction Knowledge Base. Unlike previous tools that retrieve disparate code blocks, our framework retrieves historically confirmed attack sequences, allowing the system to perform behavior-level analogical reasoning. By analyzing the semantic alignment between real-time transaction streams and historical exploit trajectories, our method can identify novel variants that share the same underlying malicious logic as past attacks, even when their syntactic features differ entirely. This provides a proactive, ``history-driven'' defense layer, effectively filling the void between static code analysis and reactive pattern matching.
